\section*{ВСТУП}

\subsection*{Постановка проблеми}

Цифрова трансформація системи вищої освіти України дедалі більше зміщується із площини організації навчального процесу у площину управління науковою діяльністю закладів вищої освіти (ЗВО). Інтеграція до європейського дослідницького простору, вимоги відкритої науки та потреба об'єктивної оцінки наукової результативності перетворюють інформаційно-комунікаційні технології (ІКТ) наукового адміністрування на критичну інфраструктуру управління освітою. Водночас обсяги наукової інформації, з якими має оперувати адміністрація ЗВО, зростають швидше за можливості традиційних, переважно ручних, процедур обліку. За даними відкритої наукометричної бази OpenAlex, світовий масив наукових праць перевищив 250 млн записів і поповнюється приблизно на 5 млн праць щороку \cite{OpenAlex2024Docs}. У таких умовах реактивна модель бібліометричної звітності, за якої заклад фіксує результати \textit{post factum} --- наприкінці звітного періоду, --- вичерпує себе й перетворюється на джерело управлінських ризиків.

Аналіз практики наукового адміністрування ЗВО дає змогу виокремити чотири взаємопов'язані проблеми, що утворюють ядро суперечності між наявними та потрібними механізмами цифрового управління наукою.

\textit{По-перше, фрагментація даних про наукову продуктивність.} Відомості про публікації одного й того самого науковця розосереджені між профілями ORCID, реферативними базами Scopus і Web of Science та відкритою базою OpenAlex. Ці джерела використовують несумісні ідентифікатори, по-різному транслітерують прізвища (наприклад, \textit{Заболотній С.} та \textit{Zabolotnii S.}), скорочують назви видань і подають метадані з різною повнотою. Унаслідок цього в типових кафедральних та факультетських звітах від 15\,\% до 25\,\% записів виявляються дублікатами однієї канонічної праці. Без коректного зведення сутностей будь-який агрегований показник --- кількість публікацій, сумарна цитованість, частка відкритого доступу --- системно спотворюється.

\textit{По-друге, надмірна трудомісткість підготовки звітності.} Формування річного звіту кафедри включає копіювання записів із кількох систем, звіряння варіантів однієї праці, усунення дублікатів і форматування переліків відповідно до вимог Міністерства освіти і науки України щодо фахових видань \cite{MONUkraine2023ProfessionalJournalsOrder}. У зовнішніх дослідженнях адміністративного навантаження науковців звітні й супровідні процедури також фіксуються як суттєва частка дослідницького часу: у FDP Faculty Burden Survey 6081 PI/Co-PI повідомили, що в середньому 42\,\% часу, відведеного на federally funded research, витрачається на адміністративні задачі, серед яких подання звітів про прогрес грантів входить до найвідчутніших навантажень \cite{Rockwell2009FDPBurden}. Тому інтервал 8--12 год для малої або середньої кафедри використано в цій роботі не як універсальний емпіричний норматив, а як консервативний калібрувальний орієнтир для подальшої прозорої моделі ручної праці. Така праця є рутинною, погано масштабованою та схильною до помилок, а її обсяг зростає з кількістю науково-педагогічних працівників і публікацій.

\textit{По-третє, реактивність моделі управління.} Чинні інформаційні процедури фіксують уже здобуті результати, але не підтримують випереджального аналізу: керівництво ЗВО не отримує своєчасних сигналів про відставання окремих підрозділів, про публікації в журналах із сумнівною репутацією чи про незадеклароване співавторство. Управлінські рішення приймаються із запізненням на цілий звітний цикл.

\textit{По-четверте, відсутність науково обґрунтованих механізмів інтеграції генеративного штучного інтелекту (GenAI) у наукове адміністрування.} Попри стрімке поширення великих мовних моделей, концептуальні моделі та відтворювані механізми їх застосування саме до завдань управління науковою діяльністю ЗВО фактично не розроблені, а наявні поодинокі рішення мають характер непрозорих евристик без верифікованої якості.

Сукупність цих проблем зумовлює \textit{актуальність} розроблення концептуальної моделі та технологічних механізмів переходу від реактивної бібліометричної звітності до проактивного управління науковою діяльністю ЗВО засобами ІКТ. Розв'язання цього завдання має водночас управлінсько-освітній та технологічний вимір, що й визначає його належність до проблематики ІКТ в управлінні освітою.

\subsection*{Аналіз останніх досліджень і публікацій}

Застосування генеративного ШІ у вищій освіті є предметом інтенсивних досліджень, проте їх тематичний фокус виявляє суттєву прогалину. Систематичний огляд Т.~Тіллманнса та співавторів показує переважання дидактичних сценаріїв GenAI у вищій освіті \cite{Tillmanns2025SystematicReviewGenAIHE}; емпіричне дослідження А.~Альшамі та співавторів деталізує сприйняття таких інструментів студентами й академічним персоналом \cite{Alshamy2025PerceptionsGenAIHE}. Натомість \textit{адміністративна} сфера, зокрема управління науковою діяльністю та наукова звітність ЗВО, лишається майже поза увагою. Саме ця незаповнена ніша визначає простір новизни нашого дослідження: ми переносимо інструментарій GenAI з дидактичної площини у площину цифрового управління наукою.

Концептуальним підґрунтям такого перенесення слугує сфера управління дослідницькою інформацією (Research Information Management, RIM) та системи поточної дослідницької інформації (Current Research Information Systems, CRIS) \cite{Bryant2017RIMDefining}. Розвиток цих систем у Європі координується ініціативами OpenAIRE та Європейською хмарою відкритої науки (EOSC), які формують вимоги до сумісності й якості дослідницьких даних \cite{OpenAIRE2023GuidelinesCRIS}, а семантичну сумісність метаданих про наукову діяльність забезпечує стандарт CERIF \cite{euroCRIS2024CERIF}. Для українських ЗВО опанування цього інструментарію є складовою євроінтеграції наукової сфери. Проте класичні CRIS зорієнтовані переважно на пасивне зберігання та обмін даними й не передбачають автономного, інтелектуального опрацювання фрагментованих джерел --- саме тут відкривається можливість для якісно нового рівня автоматизації.

Технологічною передумовою такого рівня стала поява механізмів керованої взаємодії мовних моделей із зовнішніми інструментами та даними. Протокол Model Context Protocol (MCP), оприлюднений компанією Anthropic у листопаді 2024 р. \cite{Anthropic2024MCP}, стандартизує підключення автономних агентів до зовнішніх джерел через виклик функцій (Function Calling) і дає змогу будувати агентів, здатних не лише генерувати текст, а й виконувати верифіковані операції над даними наукових профілів. Це перетворює GenAI з генератора відповідей на компонент керованого обчислювального конвеєра.

Водночас застосування мовних моделей в управлінні, де ціна помилки висока, наражається на принциповий ризик --- генерування правдоподібних, але хибних тверджень (галюцинацій). Емпіричні дослідження фіксують відчутну частоту таких помилок навіть у спеціалізованих задачах опрацювання наукової літератури \cite{Chelli2024HallucinationRates}. Звідси випливає методологічна вимога: відповідальні управлінські рішення (зокрема об'єднання записів про публікації) не можна повністю делегувати мовній моделі. Детерміновані, перевірювані процедури мають відповідати за критичні операції, а людина-експерт --- залишатися в контурі прийняття рішень (human-in-the-loop, HITL) у зонах невизначеності.

Завдання усунення дублікатів і зведення розрізнених записів до єдиної канонічної сутності є окремим випадком класичної проблеми розв'язання сутностей (Entity Resolution, ER). Її розв'язком у нашій моделі є побудова «золотого запису» (Golden Record) --- консолідованого, очищеного від дублікатів і збагаченого провенансом подання кожної праці. Як методологічну рамку для проєктування й оцінювання відповідного програмного артефакту обрано методологію проєктування на основі дизайн-науки (Design Science Research, DSR): класична праця А.~Хевнера та співавторів задає вимогу створення корисного артефакту \cite{Hevner2004DSR}, К.~Пефферс та співавтори формалізують процес DSR-оцінювання \cite{Peffers2007DSRMeth}, а Ш.~Грегор і А.~Хевнер уточнюють спосіб позиціонування внеску \cite{Gregor2013PositioningDSR}. DSR забезпечує наукову обґрунтованість запропонованих механізмів і відтворюваність їх апробації.

Отже, на перетині трьох напрямів --- перенесення GenAI з навчання в адміністрування науки, інтеграції автономних агентів через MCP у системи RIM/CRIS та контрольованого людиною розв'язання сутностей --- утворюється недостатньо досліджена проблемна область, опрацювання якої й становить мету цієї статті.

\subsection*{Мета і завдання статті}

\textbf{Мета статті} --- розробити та апробувати концептуальну модель і технологічні механізми інтеграції генеративного ШІ та автономних агентів для переходу від реактивної бібліометричної звітності до проактивного управління науковою діяльністю ЗВО засобами ІКТ, спроєктувавши їх на засадах інженерної рамки TRACE для побудови довірених агентних систем у сферах із високою ціною помилки \cite{Zabolotnii2026TRACE}.

Досягнення поставленої мети передбачає розв'язання таких завдань:

\begin{enumerate}
    \item обґрунтувати концептуальну модель проактивного управління науковою діяльністю ЗВО, що інтегрує GenAI та автономних агентів (MCP) у контур наукового адміністрування й узгоджується з підходами RIM/CRIS та стандартом CERIF;
    \item спроєктувати архітектуру системи за інженерною рамкою TRACE, рознісши компоненти на детерміноване ядро (L1), класичний ML (L2a), мовну модель-валідатор (L2b) та політику ескалації з наглядом людини (L3/L4), і обґрунтувати ощадливе застосування LLM за принципом ощадливості моделі (Model Parsimony) із кількісним індикатором --- коефіцієнтом обчислювальної ощадливості (CPR);
    \item формалізувати конвеєр розв'язання сутностей для побудови «золотого запису» над джерелами ORCID, OpenAlex і Scopus, зокрема процедури канонізації метаданих із коректним опрацюванням кириличних назв, блокування, парного зіставлення та кластеризації із захистом від хибних злиттів;
    \item визначити роль людини в контурі (HITL) як механізм пом'якшення ризику галюцинацій GenAI шляхом винесення невизначених випадків зіставлення на ручну перевірку;
    \item розробити відтворюваний синтетичний бенчмарк на узагальненому модельному середньому українському ЗВО для оцінювання якості дедуплікації та злиття записів без використання персональних даних;
    \item експериментально оцінити ефективність запропонованих механізмів за метриками якості (Precision, Recall, F1, B\textsuperscript{3} F1) у порівнянні з базовими підходами та проаналізувати чутливість результатів до частки дублікатів і ступеня спотворення метаданих;
    \item на основі прозорої аналітичної моделі оцінити потенціал скорочення трудомісткості наукової звітності та окреслити обмеження й загрози валідності запропонованого підходу.
\end{enumerate}
